% !TEX root = ../main.tex

\section{Conclusion}

This study investigates the persistent geographic (state-level) differences in U.S. stock market participation (SMP) rates, a phenomenon that standard demographic and financial characteristics cannot fully explain. Utilizing household-level panel data from the PSID spanning 1999 to 2023, paired with state-level political orientation metrics derived from FEC campaign contributions, the analysis tests whether the political environment of a state acts as a structural driver of financial behavior. The empirical tests evaluate both the initial decision (STOCKHOLD) and the allocation conditional on participation (STOCKSHARE). The most important finding is the presence of a Simpson's paradox: while individual-level shifts toward Republican preferences increase the likelihood of investing, aggregate baseline participation rates are systematically higher in Democratic-leaning states. This state-level political influence acts exclusively on the initial fixed-cost hurdle of entering the market, without affecting the conditional equity share given participation. 

By disentangling these effects, this paper contributes to the literature by bridging the gap between individual-level behavioral finance and state-level political geography to explain the persistent geographic variations in U.S. stock market participation. While a shift toward Republican preferences increases an individual's likelihood of investing,  aggregate baseline participation rates are systematically higher in Democratic-leaning states, and this state-level difference dominates. Notably, this state-level political influence acts exclusively on the initial fixed-cost hurdle of entering the market, without affecting the conditional equity share given participation. 

A major limitation of this paper is that the specific underlying mechanism or channel is yet unclear. Potential channels include corporate industry concentration, urban liquidity constraints forcing liquid investments, or cultural differences in financial literacy and information-sharing, etc., but these remain as speculative conjectures that require further causal investigation. Future research could empirically test these channels by merging additional state-level data on properties such as urbanization, industry concentration, education policies, etc. 

Another limitation is that the state-federal alignment measure's impact may have been underestimated in this methodology. The stock holdings data from PSID is biennial whereas political variables are updated every four years. Furthermore, the sample period saw frequent midterm election reversals, suggesting that alignment may have weakened or strengthened on a more short-term basis. This could have potentially diluted the alignment effects and led to the insignificant findings. However, extending the interpretation from absolute directional analysis that state politics reflect structural differences rather than proxies for individual preferences, this sounds unlikely. 